% !TEX root = ../main.tex
%-----------------------------------------------------------------------
%  The adjoint-validation result, the sequence, and where input is wanted.
%  \input by ../main.tex -- not compiled on its own.
%-----------------------------------------------------------------------
\hd{A second result: the first real check on this adjoint.}
The ensemble validates the adjoint as a by-product. The adjoint returns a
derivative at the reference point; the ensemble returns a finite difference
across members. For small perturbations the two should agree, and where they
part company marks the perturbation size beyond which the linearisation stops
holding --- which tells us whether anything built on these gradients has to be
nonlinear in \kv.

This matters more than it may sound, because \textbf{nothing else currently
checks this adjoint.} The model's built-in gradient check runs on every job
here, but it perturbs a point near the southern boundary while the cost section
sits at 26\dg N. Sensitivity there is nearly eight orders of magnitude below the
field maximum, so the check divides its own run-to-run noise by the perturbation
size and reports the result as a gradient. Both perturbed runs land on the
\emph{same side} of the base value, which a first derivative cannot produce. It
fails by roughly eight orders of magnitude, and always has.

The cause is where the check is taken, not the adjoint. Moving it into the
sensitive region is cheap and worth doing once --- but the ensemble's own
comparison is the stronger check of the same property, and would be the first
meaningful validation this configuration has had.

\hd{Sequence, and where the decision point sits.}

\smallskip
\begin{tabularx}{\textwidth}{@{}cY{1.85}Y{0.45}@{}}
\toprule
\textbf{Step} & \textbf{Content} & \textbf{Cost}\\
\midrule
1 & Build, verify against the reference result, and benchmark to confirm both
    the adjoint/forward ratio and the gradient-check saving & $<$ 1 node-h\\
2 & Generate the seven diffusivity fields and configurations & ---\\
3 & Pilot member, forward leg: check overturning drift and convective activity
  & $\sim$1.6 h\\
4 & Pilot adjoint; compare against the reference & $\sim$9 h\\
5 & \textbf{Review --- confirm the design before committing the remainder}
  & ---\\
6 & Remaining six members, run concurrently & $\sim$64 node-h\\
7 & Cross-member comparison; finite-difference validation of the gradient
  & ---\\
\bottomrule
\end{tabularx}
\smallskip

Steps 3 and 4 are one of seven runs, about $15\,\%$ of the ensemble, and would
expose a design problem before the rest is committed. The strongly perturbed
member is used for the pilot because it gives the clearest signal on both of the
things that could invalidate the design.

\hd{Where input would help most.}
Two questions change the experiment if they are answered differently. Both are
cheap to settle now and expensive to correct later, and both are better settled
before the pilot runs.

\begin{itemize}[nosep,leftmargin=1.3em]
  \item \textbf{Is ten years enough re-equilibration?} If the overturning is
        still drifting at year ten, the sensitivities partly reflect drift
        rather than a new equilibrium. The pilot member tests this directly.
  \item \textbf{Should the perturbation carry spatial structure?} The reference
        diffusivity is uniform, so as proposed every member is a constant field.
        A real ocean mixes harder over rough topography, and adding structure
        later means re-running the ensemble.
\end{itemize}
\smallskip

A third is worth flagging rather than asking about. A large implicit diffusivity
already applies wherever the column is unstable, so at the high-mixing end
explicit mixing may suppress convection --- part of the response would then be
convection triggering less often rather than a smooth dependence on \kv.
Convective activity will be diagnosed per member so the two can be told apart.

\smallskip
{\footnotesize\emph{A longer companion note covers all of this in detail, along
with a second phase that would use the ensemble as training data.}}
